% Source for the figure of HuffmanWaveletTree (sec. 22.9): the wavelet tree
% of abracadabra shaped by the canonical code a 0, b 100, c 101, d 110,
% r 111, each internal node with its items (above) and their next bits.
% Build: make figures (in ../build).
\documentclass[border=4pt]{standalone}
\usepackage[T1]{fontenc}
\usepackage{lmodern}
\usepackage{amsmath}
\usepackage{tikz}
\usepackage{lamalunga}
% tikz-ads.tex: the TikZ styles shared by the figures of the memory of the
% manual, \input by each standalone figure: a cell of a vector (cell), the
% cell that an operation touches (touched), a free cell (free), a node of a
% linked structure (node), the labels of indices (index) and of variables
% (var), a pointer (ptr), and \cells and \indices to draw a row of cells.
\usetikzlibrary{arrows.meta,positioning,calc,fit,shapes.multipart}
\definecolor{adsTouched}{RGB}{214,196,236}
\definecolor{adsFree}{RGB}{236,236,236}
\tikzset{
  ads/.style={>={Stealth[length=2mm]},font=\small,
              every node/.append style={inner sep=0pt}},
  cell/.style={draw,minimum width=7mm,minimum height=7mm,
               font=\small\ttfamily,outer sep=0pt},
  touched/.style={cell,fill=adsTouched},
  free/.style={cell,fill=adsFree,text=black!45},
  node/.style={draw,rounded corners=1pt,minimum width=8mm,
               minimum height=6mm,font=\small\ttfamily},
  index/.style={font=\scriptsize\ttfamily,text=black!60},
  var/.style={font=\small\itshape},
  ptr/.style={->,thick}
}
% \cells{name}{x,y}{list}: a row of cells, the first at (x,y), named
% name-0, name-1, ...; each item is text/style, the style one of cell,
% touched and free (e.g., {a/cell,e/touched,/free})
\newcommand{\cells}[3]{%
  \foreach \v/\s [count=\i from 0] in {#3}
    \node[\s] (#1-\i) at ($(#2)+(\i*7mm,0)$) {\v};}
% \indices{name}{n}: the indices 0 .. n-1 under the cells of name
\newcommand{\indices}[2]{%
  \foreach \i in {0,...,\numexpr#2-1\relax}
    \node[index,below=1.2mm of #1-\i] {\i};}

\begin{document}
\begin{tikzpicture}[ads,b/.style={cell,minimum width=5mm,minimum height=5.5mm},
    lf/.style={node,minimum width=6mm},sm/.style={font=\scriptsize\ttfamily,text=adsInk!60}]
  % a node: \wt{name}{x}{y}{items}{bits}
  \newcommand{\wt}[5]{%
    \foreach \c/\v [count=\i from 0] in {#4}
      {\node[b] (#1-\i) at (#2+\i*5mm,#3) {\v};
       \node[sm,above=0.6mm of #1-\i] {\c};}}
  \wt{r}{0mm}{0mm}{a/0,b/1,r/1,a/0,c/1,a/0,d/1,a/0,b/1,r/1,a/0}{}
  \wt{u}{25mm}{-17mm}{b/0,r/1,c/0,d/1,b/0,r/1}{}
  \wt{v}{10mm}{-34mm}{b/0,c/1,b/0}{}
  \wt{w}{45mm}{-34mm}{r/1,d/0,r/1}{}
  \node[lf] (la) at (2mm,-17mm) {a};
  \node[lf] (lb) at (6mm,-49mm) {b};
  \node[lf] (lc) at (24mm,-49mm) {c};
  \node[lf] (ld) at (41mm,-49mm) {d};
  \node[lf] (lr) at (59mm,-49mm) {r};
  \draw[edge] (r-0.south west) -- node[index,left=0.4mm] {0} (la.north);
  \draw[edge] (r-10.south east) -- node[index,right=0.4mm] {1} (u-3.north);
  \draw[edge] ([yshift=-3mm]u-0.south west) -- node[index,left=0.4mm] {0} ([yshift=3mm]v-1.north);
  \draw[edge] ([yshift=-3mm]u-5.south east) -- node[index,right=0.4mm] {1} ([yshift=3mm]w-1.north);
  \draw[edge] (v-0.south) -- node[index,left=0.4mm] {0} (lb.north);
  \draw[edge] (v-2.south) -- node[index,right=0.4mm] {1} (lc.north);
  \draw[edge] (w-0.south) -- node[index,left=0.4mm] {0} (ld.north);
  \draw[edge] (w-2.south) -- node[index,right=0.4mm] {1} (lr.north);
  \node[var,anchor=west] at (68mm,0) {root: the first bits of the codes};
  \node[var,anchor=west] at (68mm,-17mm) {node 1: the second bits};
  \node[var,anchor=west] at (68mm,-34mm) {nodes 10 and 11: the third bits};
\end{tikzpicture}
\end{document}
